% =====================  SECTION I  =====================
\section{Introduction}\label{sec:intro}

\subsection{Background and Motivation}\label{subsec:background}
Retail investing and digital finance have put sophisticated financial decisions in
the hands of individuals who rarely have formal training in risk management.
Artificial-intelligence tools now promise to democratize functions once reserved for
professional advisors, from automated portfolio construction to personalized
financial guidance \cite{ref1,ref2,ref8}, and machine learning has become a standard
instrument for financial forecasting and for flagging individuals at financial risk
\cite{ref3,ref7,ref18}. The same automation that makes these tools scalable also
makes them hard to see into: a retail user is routinely asked to trust a single score
with no explanation of how it was reached and no durable record of what was
recommended or why.

Three concerns have moved to the center of this literature. One is explainability: as
neural models enter financial tasks, methods that attribute a decision to its inputs
have become essential to using them responsibly \cite{ref10,ref12,ref21}. A second is
trust and reliance---behavioral studies of AI financial advisors find that disclosing
a system's limitations, rather than presenting it as infallible, produces more
calibrated user reliance \cite{ref11}. The third, verifiability, follows from the
stakes involved: because financial recommendations carry consequences, tamper-evident
audit trails, often built on blockchain-style hashing, let a recommendation be checked
afterward against what was actually issued \cite{ref13,ref16,ref22}. Cryptocurrency
allocation sharpens all three, since the assets are volatile and the
decentralized-finance setting adds liquidation and stablecoin-stability risks
\cite{ref5,ref6}.

\subsection{Problem Statement and Existing Limitations}\label{subsec:problem}
Despite substantial progress on each concern individually, they are almost always
addressed in isolation. Robo-advisory and portfolio systems optimize allocations but
seldom expose an interpretable risk rationale \cite{ref8,ref12}; distress- and
risk-prediction models classify users but stop short of translating a prediction into
an actionable, safety-constrained recommendation \cite{ref18,ref19}; explainability
research develops attribution methods largely in the abstract \cite{ref10,ref21}; and
blockchain-audit research secures records without connecting them to a live decision
process \cite{ref13,ref14,ref15}. A retail user is thus served by a patchwork of
tools, none of which delivers a recommendation that is simultaneously data-driven,
interpretable, conservative under uncertainty, and independently verifiable. A further
and subtler limitation pervades applied financial machine learning: reported
accuracies are frequently inflated by information leakage between how labels are
constructed and what features a model is given, and by evaluation protocols that do
not account for repeated entities \cite{ref4}. Systems are seldom candid about the gap
between a headline metric and the quality of the decision a user ultimately receives.

\subsection{Research Gap}\label{subsec:gap}
The gap this work addresses is therefore integrative and methodological rather than
algorithmic. No single, inspectable system in the reviewed literature carries a user's
raw financial inputs through interpretable feature construction, unsupervised risk
labeling, a leakage-aware classifier, an explicit safety layer, per-decision
explanation, risk-aligned allocation, and a verifiable audit record---while being
explicit about where information leaks, which component actually determines the
outcome, and what the evidence does and does not support. Closing this gap requires not
a new estimator but a coherent architecture with enforceable properties and an honest
evaluation of them.

\subsection{Objectives and Research Questions}\label{subsec:rqs}
The objective of this work is to design, implement, and critically evaluate an
end-to-end pipeline that integrates the five concerns above into a single
decision-support tool, and to hold that tool to a transparent evidential standard.
This objective is refined into six research questions:

\begin{itemize}
\item[\textbf{RQ1.}] Can a user's financial risk be assessed from a small set of raw,
  deployable inputs while the ratios used to construct the risk labels are withheld
  from the classifier, so that leakage can be reasoned about explicitly?
\item[\textbf{RQ2.}] Do unsupervised clusters over interpretable financial ratios
  correspond to coherent, monotonically ordered Low/Medium/High risk bands suitable as
  proxy labels?
\item[\textbf{RQ3.}] Can a supervised classifier reproduce those proxy labels from the
  raw inputs alone with useful accuracy, and is a neural model warranted over simpler
  baselines?
\item[\textbf{RQ4.}] How can a learned signal, an unsupervised signal, and a
  deterministic rule-based signal be reconciled so that the system's failure mode is to
  over-warn rather than under-warn?
\item[\textbf{RQ5.}] Can a single explanation strategy provide per-decision,
  input-level attributions across different model types without retraining?
\item[\textbf{RQ6.}] Can the reconciled risk band be mapped to a risk-aligned
  cryptocurrency allocation and logged as a tamper-evident, independently verifiable
  record?
\end{itemize}

\subsection{Contributions}\label{subsec:contributions}
Against these questions, the paper makes the following contributions:

\begin{itemize}
\item[\textbf{C1)}] An integrated, end-to-end pipeline that unifies feature
  engineering, unsupervised proxy labeling, leakage-aware classification, conservative
  reconciliation, per-decision explanation, risk-aligned allocation, deterministic
  simulation, and cryptographic audit logging in one inspectable system, surfaced
  through a working application.
\item[\textbf{C2)}] An explicit feature contract that separates the ratios used for
  clustering from the raw inputs given to the classifier, together with a measured
  analysis of the residual arithmetic leakage that column-level separation does not
  remove.
\item[\textbf{C3)}] A conservative guardrail and a most-cautious reconciliation rule
  that provably cannot lower risk, accompanied by a quantitative demonstration that
  this deterministic layer---not the learned model---determines the outcome for most
  users.
\item[\textbf{C4)}] A model-adaptive, occlusion-based explanation mechanism that yields
  per-input influence scores for any probabilistic classifier in the pipeline without
  retraining.
\item[\textbf{C5)}] A risk-aligned allocation and a tamper-evident SHA-256 audit
  record, integrated so that every issued recommendation is both explainable and
  independently checkable.
\item[\textbf{C6)}] A candid evaluation methodology in which label leakage, guardrail
  dominance, and a statistically non-significant model-selection result are measured
  and reported rather than assumed away---offered as a template for honest appraisal of
  integrated financial AI systems.
\end{itemize}

\subsection{Paper Organization}\label{subsec:organization}
The remainder of the paper is organized as follows. Section~\ref{sec:related} reviews
related work thematically and positions the contribution against it.
Section~\ref{sec:method} presents the proposed framework, the dataset, the
mathematical formulation, and each pipeline component. Section~\ref{sec:results}
describes the experimental setup and reports results, including the clustering,
classification, comparative, guardrail, simulation, and auditability findings.
Section~\ref{sec:discussion} discusses the implications, states the limitations and
threats to validity in full, and separates the engineering contribution from the
research contribution. Section~\ref{sec:conclusion} concludes and outlines future work.
